% --- Archon physics formalization source begin ---
% archon:physics
% archon:covers PhyXMiniProblems/problem_phyx_mini_0008.lean
% archon:source-report reports/phyx_mini/problem_phyx_mini_0008.source.json

\chapter{Physics problem phyx\_mini\_0008}
\label{ch:PhyXMiniProblems_problem_phyx_mini_0008}

\paragraph{Problem source.}
\#\# Physical scenario

The index of refraction for violet light in silica flint glass is 1.66, and that for red light is 1.62.

\#\# Question

What is the angular spread of visible light passing through a prism of apex angle \textbackslash{}( 60.0\textasciicircum{}\textbackslash{}circ \textbackslash{}) if the angle of incidence is \textbackslash{}( 50.0\textasciicircum{}\textbackslash{}circ \textbackslash{})? See figure.

\#\# Answer choices

- A: A: \textbackslash{}( 4.37\textasciicircum{}\textbackslash{}circ \textbackslash{})

- B: B: \textbackslash{}( 8.74\textasciicircum{}\textbackslash{}circ \textbackslash{})

- C: C: \textbackslash{}( 4.61\textasciicircum{}\textbackslash{}circ \textbackslash{})

- D: D: \textbackslash{}( 2.11\textasciicircum{}\textbackslash{}circ \textbackslash{})

\#\# Figure caption (auxiliary; use the image as primary evidence)

The image depicts a visual representation of light dispersion through a prism. Key elements include:

- A triangular prism on the left through which a beam of "Visible light" enters from the left side.
- As the light passes through the prism, it is dispersed into a spectrum of colors.
- This spectrum is shown exiting the prism on the right, spreading out into different colored rays labeled as R (red), O (orange), Y (yellow), G (green), B (blue), and V (violet).
- The spectrum of colors is projected onto a "Screen," demonstrating the dispersion effect.
- Two labels indicate specific angles: "Deviation of red light" and "Angular spread," illustrating the angles at which the light is bent and spread after passing through the prism.

\#\# Dataset metadata

Geometrical Optics; Physical Model Grounding Reasoning, Spatial Relation Reasoning

\paragraph{Recorded answer/context.}
C: C: \textbackslash{}( 4.61\textasciicircum{}\textbackslash{}circ \textbackslash{})

\paragraph{Figure/image path.}
/root/proposal\_for\_physic/science-mango/phyx\_mini\_run/phyx\_data/test\_image/8.png

\paragraph{Formalization target.}
create a compiling Lean file with sorry bodies at `PhyXMiniProblems/problem\_phyx\_mini\_0008.lean`. The Lean declarations must preserve the physical quantities, dimensions or dimensional roles, figure labels, governing-law hypotheses, and final relation expressed by this problem.
Use Mathlib/Physlib names found through LeanExplore where available. If a physics API is missing, introduce faithful local abstractions rather than scalar placeholder aliases.

\begin{theorem}[Physics formalization target]
\label{thm:physics:phyx_mini_0008:target}
\lean{PhyXMiniProblems.Problem0008.angularSpread_is_answer_C}
\uses{def:physics:phyx-mini-0008:phyxminiproblems-problem0008-degreestoradians, def:physics:phyx-mini-0008:phyxminiproblems-problem0008-prismsetup, def:physics:phyx-mini-0008:phyxminiproblems-problem0008-prismrayangles, def:physics:phyx-mini-0008:phyxminiproblems-problem0008-satisfiesprismraylaws, def:physics:phyx-mini-0008:phyxminiproblems-problem0008-angularspreadradians, def:physics:phyx-mini-0008:phyxminiproblems-problem0008-dispersionfigurereadout, def:physics:phyx-mini-0008:phyxminiproblems-problem0008-answerangledegrees}
For a $60^\circ$ silica-flint prism in air at $50^\circ$ incidence, with
dimensionless indices $n_R=1.62$ and $n_V=1.66$, impose the two Snell laws,
the triangular-prism angle relation, the principal physical branches, and the
usual deviation relation separately for red and violet rays. The
violet-minus-red deviation is approximately $4.6118^\circ$ and therefore
rounds to $4.61^\circ$, answer C. Colors, media, and ray states must remain
typed physical roles, with real numbers used only for explicit angle and
index readouts.
\end{theorem}
\begin{proof}
For each color with index $n$, first-interface Snell refraction gives
$r_1=\arcsin(\sin 50^\circ/n)$. Prism geometry gives
$r_2=60^\circ-r_1$, and the second interface gives
$e=\arcsin(n\sin r_2)$. The deviation is
$\delta=50^\circ+e-60^\circ$. Evaluating this expression at $n_V$ and
$n_R$ and subtracting gives about $4.6118^\circ$, which lies within
$0.005^\circ$ of $4.61^\circ$ and is not equally close to another choice.
\end{proof}

% --- Archon declaration topology begin ---
\section{Declaration topology}
\begin{definition}[Visible Color]
  \label{def:physics:phyx-mini-0008:phyxminiproblems-problem0008-visiblecolor}
  \lean{PhyXMiniProblems.Problem0008.VisibleColor}
  The six outgoing ray labels shown on the screen, ordered from red to violet
\end{definition}
\begin{proof}
  This is the declared model definition of Visible Color.
\end{proof}
\begin{definition}[Optical Medium]
  \label{def:physics:phyx-mini-0008:phyxminiproblems-problem0008-opticalmedium}
  \lean{PhyXMiniProblems.Problem0008.OpticalMedium}
  The two media crossed by each ray in the prism experiment
\end{definition}
\begin{proof}
  This is the declared model definition of Optical Medium.
\end{proof}
\begin{definition}[degrees To Radians]
  \label{def:physics:phyx-mini-0008:phyxminiproblems-problem0008-degreestoradians}
  \lean{PhyXMiniProblems.Problem0008.degreesToRadians}
  Convert a scalar angle read in degrees to its value in radians
\end{definition}
\begin{proof}
  This is the declared model definition of degrees To Radians.
\end{proof}
\begin{definition}[Prism Setup]
  \label{def:physics:phyx-mini-0008:phyxminiproblems-problem0008-prismsetup}
  \lean{PhyXMiniProblems.Problem0008.PrismSetup}
  \uses{def:physics:phyx-mini-0008:phyxminiproblems-problem0008-visiblecolor, def:physics:phyx-mini-0008:phyxminiproblems-problem0008-opticalmedium}
  Experimental data shared by every color: a color-dependent dimensionless refractive index, the incidence angle, and the triangular prism's apex angle
\end{definition}
\begin{proof}
  This is the declared model definition of Prism Setup.
\end{proof}
\begin{definition}[Prism Ray Angles]
  \label{def:physics:phyx-mini-0008:phyxminiproblems-problem0008-prismrayangles}
  \lean{PhyXMiniProblems.Problem0008.PrismRayAngles}
  \uses{def:physics:phyx-mini-0008:phyxminiproblems-problem0008-visiblecolor}
  The four directed angle readouts needed to trace one colored ray through the two faces of the prism
\end{definition}
\begin{proof}
  This is the declared model definition of Prism Ray Angles.
\end{proof}
\begin{definition}[Satisfies Prism Ray Laws]
  \label{def:physics:phyx-mini-0008:phyxminiproblems-problem0008-satisfiesprismraylaws}
  \lean{PhyXMiniProblems.Problem0008.SatisfiesPrismRayLaws}
  \uses{def:physics:phyx-mini-0008:phyxminiproblems-problem0008-visiblecolor, def:physics:phyx-mini-0008:phyxminiproblems-problem0008-prismsetup, def:physics:phyx-mini-0008:phyxminiproblems-problem0008-prismrayangles}
  The geometrical-optics laws for one color. The interval conditions select the physical, principal-angle branch of Snell's law
\end{definition}
\begin{proof}
  This is the declared model definition of Satisfies Prism Ray Laws.
\end{proof}
\begin{definition}[angular Spread Radians]
  \label{def:physics:phyx-mini-0008:phyxminiproblems-problem0008-angularspreadradians}
  \lean{PhyXMiniProblems.Problem0008.angularSpreadRadians}
  \uses{def:physics:phyx-mini-0008:phyxminiproblems-problem0008-prismrayangles}
  The angular separation from the red ray to the violet ray in the figure
\end{definition}
\begin{proof}
  This is the declared model definition of angular Spread Radians.
\end{proof}
\begin{definition}[Dispersion Figure Readout]
  \label{def:physics:phyx-mini-0008:phyxminiproblems-problem0008-dispersionfigurereadout}
  \lean{PhyXMiniProblems.Problem0008.DispersionFigureReadout}
  \uses{def:physics:phyx-mini-0008:phyxminiproblems-problem0008-angularspreadradians}
  The two angle labels explicitly drawn in the source figure. The colored rays terminate on the screen in the order represented by VisibleColor above
\end{definition}
\begin{proof}
  This is the declared model definition of Dispersion Figure Readout.
\end{proof}
\begin{definition}[Answer Choice]
  \label{def:physics:phyx-mini-0008:phyxminiproblems-problem0008-answerchoice}
  \lean{PhyXMiniProblems.Problem0008.AnswerChoice}
  Multiple-choice labels from the problem statement
\end{definition}
\begin{proof}
  This is the declared model definition of Answer Choice.
\end{proof}
\begin{definition}[answer Angle Degrees]
  \label{def:physics:phyx-mini-0008:phyxminiproblems-problem0008-answerangledegrees}
  \lean{PhyXMiniProblems.Problem0008.answerAngleDegrees}
  \uses{def:physics:phyx-mini-0008:phyxminiproblems-problem0008-answerchoice}
  The displayed degree value attached to each answer choice
\end{definition}
\begin{proof}
  This is the declared model definition of answer Angle Degrees.
\end{proof}
% --- Archon declaration topology end ---
% --- Archon physics formalization source end ---
